%! Inverse Trigonometric Functions — Unit 3, Lesson 3.9 — Warm-Up (Answer Key)
\documentclass[10pt]{article}
\usepackage{precalculus-article}
\usepackage{precalculus-key}   % pulls in -boxes; do NOT also load -boxes
\newcommand{\TallMath}[1]{$\displaystyle #1\rule[-1.4em]{0pt}{3.2em}$}

\begin{document}

\pageheader{Unit 3, Lesson 3.9}{Warm-Up --- Answer Key}

\vspace{0.12in}

\begin{notesbox}{Warm-Up: Prerequisite Review --- Answer Key}

\textbf{1.} Evaluate each trigonometric expression.

\renewcommand{\arraystretch}{2.2}
\begin{tabularx}{\linewidth}{@{} X X X X @{}}
  \textbf{(a)} $\sin\!\left(\dfrac{\pi}{6}\right) = $ \ans{$\dfrac{1}{2}$}
  &
  \textbf{(b)} $\cos\!\left(\dfrac{\pi}{3}\right) = $ \ans{$\dfrac{1}{2}$}
  &
  \textbf{(c)} $\tan\!\left(\dfrac{\pi}{4}\right) = $ \ans{$1$}
  &
  \textbf{(d)} $\cos\!\left(\dfrac{5\pi}{6}\right) = $ \ans{$-\dfrac{\sqrt{3}}{2}$}
\end{tabularx}

\vspace{0.10in}

\textbf{2.} The table below shows some values of a function $f$.

\begin{center}
\renewcommand{\arraystretch}{1.6}
\begin{tabular}{|c|c|c|c|c|c|}
\hline
$x$ & $0$ & $\dfrac{\pi}{6}$ & $\dfrac{\pi}{2}$ & $\dfrac{5\pi}{6}$ & $\pi$ \\
\hline
$f(x)$ & $0$ & $\tfrac{1}{2}$ & $1$ & $\tfrac{1}{2}$ & $0$ \\
\hline
\end{tabular}
\end{center}

\vspace{0.06in}
\textbf{(a)} Is $f$ a one-to-one function?  \ans{No}

\vspace{0.06in}
\textbf{(b)} Explain.

\ans{The output $\tfrac{1}{2}$ corresponds to two different inputs
($\tfrac{\pi}{6}$ and $\tfrac{5\pi}{6}$), so the function fails the
one-to-one condition.}

\vspace{0.06in}

\textbf{3.} Domain and range of $f(x) = \sin x$ (unrestricted).

\medskip
Domain: \ans{all real numbers} \qquad
Range: \ans{$[-1,\, 1]$}

\vspace{0.06in}

\textbf{4.} For a function to have an inverse, it must be
\ans{one-to-one (injective)}.
Graphical test: \ans{the Horizontal Line Test --- every horizontal line
intersects the graph at most once.}
\end{notesbox}

\end{document}
